\section{Framework}
\label{sec:framework}

This section develops a framework that predicts, from a language's grammar alone, how much morphology-aware tokenization lowers the parameter count a transformer needs to reach a fixed capability threshold. The framework is built in six moves. \S\ref{sec:framework-ssc} defines the \textbf{Structural Synthesis Ceiling}, a structural upper bound on the grammatical information a single word form can carry, and sets it apart from Joseph Greenberg's corpus-averaged Index of Synthesis. \S\ref{sec:framework-H} gives the information-theoretic quantity $H(L)$, the grammatical Shannon entropy per word form. \S\ref{sec:framework-rho} introduces the \textbf{structural recoverability coefficient} $\rho(L)$, which captures how much of $H(L)$ is recoverable by a regular surface-only parser. \S\ref{sec:framework-A1} states the assumption (A1) that ties grammatical information to model capacity, and names the \emph{implicit grammar cost} a BPE-tokenized transformer must pay. \S\ref{sec:framework-engine} describes the per-language engine architecture (a morphology layer plus a grammar layer, together with the categorical input streams they emit) through which $B(L)$ becomes a concrete embedding table. \S\ref{sec:framework-rebate} combines these into a closed-form prediction for the \textbf{parameter rebate} $\Delta N(L)$ that morphology-aligned tokenization delivers relative to BPE, and spells out the corresponding \textbf{capability-vs-scale curve shift}. The empirical status of this prediction, with ordinal confirmation and magnitude divergence, is the subject of \S\ref{sec:results}.

\subsection{The Structural Synthesis Ceiling}
\label{sec:framework-ssc}

Consider a language $L$ whose morphology realises grammatical features on word forms through any combination of inflection, derivation, affixation, and non-concatenative processes. For each grammatical category $C$ encoded in $L$'s morphology (tense, aspect, mood, person, number, gender, case, definiteness, possession, polarity, and so on, the inventory varying from language to language), let $V(C)$ be the set of values that category admits. Let $B(L)$ be the set of distinct grammatical feature bundles that $L$'s morphology can realise simultaneously on a single word form, once syncretism, defectiveness, and paradigm gaps have been taken into account. $B(L)$ is, in effect, the realised subset of the Cartesian product $\prod_C V(C)$, the structural ``address space'' of a word form in $L$.

We define the \textbf{Structural Synthesis Ceiling} (SSC) of $L$ as the cardinality of $B(L)$, and its logarithm in bits:

\begin{equation}
\mathrm{SSC}(L) = |B(L)|, \qquad H_{\max}(L) = \log_2 |B(L)|.
\label{eq:ssc}
\end{equation}

$H_{\max}(L)$ is the maximum grammatical information, in bits, that a single word form can carry, independent of any corpus, tokenizer, or training run. It can be read off a language's grammar description alone, and in our implementation it is enumerated directly from the feature-bundle registry of each language's grammar engine.

SSC sits within a lineage of typological indices, but differs from its best-known predecessor. \citet{greenberg1960} defined the \textbf{Index of Synthesis} as the average number of morphemes per word, computed over a corpus. The Index of Synthesis is a \emph{descriptive} typological measure, telling us how many morphemes are strung together in a representative sample of text. SSC is instead a \emph{structural upper bound}, telling us how many distinct grammatical specifications the morphology can in principle express on a single word, regardless of how often each is realised in practice. The Index of Synthesis is empirical and averaged; SSC is combinatorial and peaked. For our purposes, predicting how much grammatical information a tokenizer can pre-encode, the structural ceiling is the relevant quantity.

For the three languages this paper tests, SSC spans the morphological typology we set out to cover. English, an analytic language whose grammar is carried largely by word order, has $|B(L)| = 23$ ($H_{\max} \approx 4.52$ bits). Turkish, an agglutinative language whose word forms stack suffixes in a strict slot order, has $|B(L)| = 63$ ($H_{\max} \approx 5.98$ bits). Arabic, a root-and-pattern templatic language whose grammar arises from interleaving consonantal roots (mostly triliteral, with a smaller class of quadriliteral roots and a handful of longer ones) with vowel patterns (\ar{أَوْزَان}, \emph{awz\=an}), has $|B(L)| = 270$ ($H_{\max} \approx 8.08$ bits), more than four times Turkish's structural bundle space despite Turkish's combinatorially spread agglutinative slot system.

A note on terminology. Throughout this paper $|B(L)|$ denotes the \emph{structural} bundle-space size as enumerated by the grammar engine's feature-bundle registry, that is, the number of distinct grammatical specifications the morphology can in principle realise on a single word form. This is a property of the grammar, not of any corpus. The empirical-distribution support size, namely the number of distinct \texttt{feature\_bundle\_str()} outputs observed in a given corpus sample, can deviate from $|B(L)|$ in both directions. Arabic's 270 structural bundles surface as roughly 697 distinct strings in a 20{,}000-sentence Wikipedia sample, because the engine serialises tag combinations as distinct strings, while Turkish's 63 structural bundles expand to roughly 1{,}837 observed strings through the combinatorial spread of slot values. These observed counts enter the $H(L)$ computation in \S\ref{sec:framework-H}, which is distribution-weighted and so invariant to such expansion, but the canonical framework quantities remain the structural $|B(L)|$ and its logarithm $H_{\max}(L)$.

\subsection{Grammatical information content $H(L)$}
\label{sec:framework-H}

$H_{\max}(L)$ is only an upper bound. The grammatical information actually carried by an average word form depends on how bundles are distributed in naturally occurring text. Let $p(b)$ denote the empirical probability, over a corpus of $L$, that a given word form carries feature bundle $b \in B(L)$. The \textbf{grammatical information content} of $L$, in bits per word form, is the Shannon entropy over this distribution:

\begin{equation}
H(L) = -\sum_{b \in B(L)} p(b) \, \log_2 p(b).
\label{eq:H}
\end{equation}

$H(L) \le H_{\max}(L)$, with equality only when every bundle is equiprobable. In practice the realised distribution is highly skewed. A handful of common bundles (third-person singular present in English, nominative singular in Turkish, Form~I active perfective in Arabic) dominate, and $H(L)$ falls well below the ceiling. On our 20{,}000-sentence Wikipedia samples we observe $H(L) = 1.43$ bits for English, $4.97$ bits for Turkish, and $5.44$ bits for Arabic.

The ordering of $H(L)$, namely English $\ll$ Turkish $\approx$ Arabic, already departs from what the SSC alone would predict. Arabic's ceiling sits 35\,\% above Turkish's in bits, yet the corpus-weighted entropies are nearly equal. The difference is absorbed by the structural recoverability coefficient introduced next.

\subsection{Structural recoverability $\rho(L)$}
\label{sec:framework-rho}

$H(L)$ measures the grammatical information a word form carries. It does not measure how much of that information is accessible to a surface-only parser, that is, a parser that sees the word's surface string and nothing else, with no access to a lexicon and no context. The distinction matters because tokenizers, whether BPE or morphology-aligned, work on surface forms.

We capture this distinction with a single coefficient. Let $s \in \Sigma$ range over surface forms and $b \in B(L)$ over feature bundles. The \textbf{conditional entropy of the bundle given the surface form} is:

\begin{equation}
H(b \mid s) = \sum_{s} p(s) \cdot H(b \mid s = s),
\label{eq:Hcond}
\end{equation}

where $H(b \mid s = s)$ is the Shannon entropy of the bundle distribution conditional on the specific surface form $s$, estimated from corpus-observed $(s, b)$ pairs. We define the \textbf{structural recoverability coefficient} of $L$ as:

\begin{equation}
\rho(L) = 1 - \frac{H(b \mid s)}{H(L)}.
\label{eq:rho}
\end{equation}

Equation~\ref{eq:rho} reads off directly. $H(b \mid s)$ measures the grammatical uncertainty that remains once the surface form has been observed. When the morphology is one-to-one with its surface realisation, with each bundle corresponding to a unique surface string as in an idealised agglutinative language, $H(b \mid s) = 0$ and $\rho = 1$. When the surface form carries no information at all about the bundle, the degenerate case of complete ambiguity, $H(b \mid s) = H(L)$ and $\rho = 0$. Real languages fall in between.

Two technical points deserve mention. First, the conditional entropy~\ref{eq:Hcond} is conditioned on the surface string itself, rather than on a ``surface signature'' produced by stripping the lexical content. For concatenative morphologies such as English and Turkish, a lemma-stripped signature differs from the full surface only in lexical root identity, which in principle carries no grammatical information. For templatic morphologies such as Arabic, the grammar engine's explicit \texttt{template} field, a categorical label naming the morphological pattern (for instance \texttt{VERB\_TRILATERAL\_BARE} or \texttt{NOM\_DERIVED}), plays this role directly. Second, $\rho(L)$ is a property of the language's grammar together with its corpus distribution, not of any particular tokenizer. Both BPE and morphology-aligned tokenizers, as surface-only systems, are bounded above in the grammatical information they can recover by $H_{\mathrm{eff}}(L) = \rho(L) \cdot H(L)$.

On the 20{,}000-sentence samples we find $\rho(\text{English}) \approx 0.89$, $\rho(\text{Turkish}) \approx 0.85$, and $\rho(\text{Arabic}) \approx 0.59$. The ordering reflects the typological character of the three systems. English and Turkish, both fundamentally concatenative, expose most of their grammatical structure on the surface. Arabic's templatic fusion hides roughly 40\,\% of the grammatical signal behind ambiguities that lexical context or full analysis would be required to resolve. The absolute values are signature-proxy estimates. A more refined estimator, for instance the mutual information between surface suffix $n$-grams and bundles, is left to follow-up work.

\subsection{Implicit grammar cost and Assumption A1}
\label{sec:framework-A1}

A transformer trained on byte-pair tokens receives no structural information about morphology at its input. Grammatical structure (which bundles go with which lexical material, which inflections follow which roots, which surface patterns realise which feature combinations) has to be learned from co-occurrence statistics in the training data. That learning consumes model capacity.

Let $C_{\mathrm{implicit}}(L)$ denote the model capacity, in bits, that a BPE-trained transformer must allocate to reconstructing $H_{\mathrm{eff}}(L)$ from the distribution, at a capability threshold to be specified. Under morphology-aligned tokenization, the same grammatical information is instead handed to the model as an explicit categorical input: a bundle-embedding table indexed by $b \in B(L)$, summed into the model's input representation alongside the root embedding. The capacity cost of the explicit path is just the storage of the embedding table:

\begin{equation}
C_{\mathrm{explicit}}(L) = |B(L)| \cdot d_{\mathrm{embed}} \cdot (\text{bits per weight}).
\label{eq:Cexplicit}
\end{equation}

For the Phase~1 configuration ($d_{\mathrm{embed}} = 128$, 32-bit precision, $|B(L)|$ at most about $1{,}800$ on our Turkish 20k sample), $C_{\mathrm{explicit}}$ is on the order of 10\,MB, three or four orders of magnitude below a 100M-parameter model. For the purposes of the rebate argument we therefore treat $C_{\mathrm{explicit}}(L)$ as zero.

The harder quantity to pin down is $C_{\mathrm{implicit}}(L)$. A precise derivation from first principles would require a theory of how transformers allocate parameters to distributional pattern-learning, and no such theory exists in tractable form. We therefore promote a functional form to a stated assumption, to be examined empirically in \S\ref{sec:results}.

\paragraph{Assumption A1 (dominant-term linearity).}
The capacity a BPE-trained transformer allocates to implicitly reconstructing grammatical structure is, to first order, linear in the effective grammatical information content $\rho(L) \cdot H(L)$:
\begin{equation}
C_{\mathrm{implicit}}(L) = \alpha \cdot \rho(L) \cdot H(L) + O(\text{lower-order terms}).
\label{eq:A1}
\end{equation}

The coefficient $\alpha$ absorbs factors we take to be language-invariant at fixed architecture, data scale, and tokenizer family: the inefficiency of BPE fragments as morphological units, the effective vocabulary over which the grammatical association has to be learned, and the bits-per-parameter capacity of the underlying transformer \citep{allenzhu2024}. A1 is the simplest non-trivial hypothesis about how grammatical information translates into model capacity. It is not a theorem, it is the prediction the framework commits to, and what gets tested below.

\subsection{Engine architecture: the morphology layer, the grammar layer, and the input streams}
\label{sec:framework-engine}

The previous subsections treated $|B(L)|$, $H(L)$, and $\rho(L)$ as abstract quantities derivable from a language's grammar. They are also concrete outputs of a piece of software, the per-language \emph{morphology engine}, whose architecture we describe here. The engine matters to the framework for two reasons. It is the operational source of the bundle inventory $B(L)$ that anchors Equation~\ref{eq:ssc}, and it determines the shape of the input that the morphology-aligned transformer receives, which in turn fixes what is meant by ``explicit categorical input'' in \S\ref{sec:framework-A1}.

\paragraph{Two layers, one composition.}
Each language is served by a two-layer engine. The lower layer, the \emph{morphology engine}, operates per word. For a given surface form it emits a \texttt{TokenInfo} record containing the surface string, the root or stem, a feature-bundle (part of speech together with the grammatical tags the morphology realises on that word), and a small number of language-specific extras (templatic class for Arabic, radical class for Mandarin, and so on). The upper layer, the \emph{grammar layer}, operates per sentence. It does two things. First, it uses neighbour context to resolve part-of-speech ambiguities that no per-word analyser can settle on its own (the English \emph{run} verb-or-noun question, the Turkish verbal-noun-versus-finite-verb question, the Arabic \ar{مَدْرَسَة} school-or-teacher-feminine question). Second, it validates the sentence against a small set of permitted structural templates encoded as Python predicates, with rejection messages written in the host language's own grammatical terminology (\ar{الْفَاعِل} for Arabic subject position, \emph{fiil-sonda} for Turkish verb-final order, \zh{把构式} for the Mandarin \emph{b\v{a}}-construction, and so on). The morphology proposes; the grammar disposes. The composition matters: bundle counts $|B(L)|$ reported in \S\ref{sec:framework-ssc} are the inventories surfaced after grammar-layer disambiguation, not the larger noisy union the morphology alone would emit.

\paragraph{Four engines, distinct stream counts.}
The framework is instantiated across four languages, each chosen as a representative of a different morphological type. Each engine emits a fixed number of categorical streams per token, and that stream count determines the input shape the transformer downstream consumes.

\begin{itemize}
\item \textbf{Mandarin (ZH, isolating with radical-class semantics).} Two streams per token: the character itself, and its Kangxi radical class. The radical inventory is the canonical 214-class Kangxi system, sourced from the Unihan database's \texttt{kRSKangXi} field via \texttt{configs/zh\_radicals.json}.
\item \textbf{English (EN, analytic).} Two streams per token: surface form and feature-bundle. Morphological analysis is a recursive multi-pass decomposition (inflection pass, suffix pass, prefix pass), with explicit guards (\texttt{NO\_MENT\_STRIP}, \texttt{NO\_PREFIX\_PEEL\_BASES}, and so on) that prevent over-stripping Latin-fused words for which surface affixes are not synchronically productive.
\item \textbf{Turkish (TR, agglutinative).} Two streams per token: surface form and feature-bundle. Strict-slot suffix peeling with vowel-harmony normalisation, applied holistically across verbal and nominal paradigms, with deterministic tie-breaks when more than one slot sequence could in principle parse a given surface string.
\item \textbf{Arabic (AR, root-and-pattern templatic).} \emph{Four} streams per token: surface form, feature-bundle, root (a triliteral, quadriliteral, or longer consonantal skeleton, drawn from a 7{,}142-entry registry in \texttt{ar\_roots.json}), and a \emph{wazn-class semantic role} drawn from the 185-value inventory in \texttt{ar\_templates.json}'s \texttt{semantic\_role} field.
\end{itemize}

\paragraph{The wazn-class taxonomy in full.}
The Arabic wazn-class stream deserves its own description, because it is the framework's principal departure from a purely concatenative view of morphology. A \emph{wazn} (\ar{وَزْن}, \emph{wazn}, plural \ar{أَوْزَان} \emph{awz\=an}) is a morphological pattern: a template of vowels and consonant slots into which a root is interleaved to yield a word. The engine surfaces 185 distinct \texttt{semantic\_role} values, partitioned into three broad groups. The verbal forms run from \texttt{FORM\_I\_BARE} through \texttt{FORM\_XII\_INTENSIVE\_HABITUAL}, covering the classical Form~I through Form~X system together with the rarer Forms~XI and~XII. The derived nominal classes name the standard nominal patterns of Arabic morphology: \texttt{AGENT\_NOUN} (\ar{اسْم الْفَاعِل}), \texttt{PATIENT\_NOUN} (\ar{اسْم الْمَفْعُول}), \texttt{INSTRUMENT\_NOUN} (\ar{اسْم الْآلَة}), \texttt{PLACE\_TIME} (\ar{اسْم الزَّمَان وَالْمَكَان}), \texttt{MASDAR\_MIMI}, \texttt{MASDAR\_MARRA}, \texttt{MASDAR\_HAYAA}, \texttt{SIFA\_MUSHABBAHA} (\ar{الصِّفَة الْمُشَبَّهَة}), \texttt{MUBALAGHAH} (\ar{صِيغَة الْمُبَالَغَة}), \texttt{NISBA} (\ar{النِّسْبَة}), \texttt{ELATIVE} (\ar{اسْم التَّفْضِيل}), and others. The remainder is a long tail of weakness-specific variants that name the same semantic role with the root's weak-letter behaviour spelled out (\texttt{ACTION\_AJWAF\_WAW}, \texttt{VERBAL\_NOUN\_HOLLOW\_FUWAAL}, and so on), because a hollow root's \emph{wazn} and a sound root's \emph{wazn} are not the same surface object and merging them would obscure the very signal the stream is meant to provide. The total, 185 distinct values, is the inventory the embedding table indexes.

\paragraph{Input embedding architecture.}
The transformer downstream is identical across regimes: same depth, same width, same attention configuration, same head count. Parameter parity holds everywhere except at the input embedding layer. Per token, the model receives $N$ categorical streams (one for baseline BPE, two for ZH, EN, and TR under morphology-aligned tokenization, four for AR), and each stream owns its own learned embedding table. The streams are combined by summation at the input:

\begin{equation}
x = \sum_{s \in \mathrm{streams}} E_s[\mathrm{id}_s],
\label{eq:embedsum}
\end{equation}

so that for Arabic specifically the input representation reads

\[
x = E_{\mathrm{tok}}[\mathrm{tok}] + E_{\mathrm{feat}}[\mathrm{feat}] + E_{\mathrm{root}}[\mathrm{root}] + E_{\mathrm{wazn}}[\mathrm{wazn}],
\]

with each table sharing the model dimension $d_{\mathrm{embed}}$. This is the operational meaning of ``explicit categorical input'' in \S\ref{sec:framework-A1}: the grammar that a BPE-trained transformer must reconstruct distributionally is instead handed to the morphology-aligned transformer as a vector lookup, summed into the same representation the attention stack reads.

\subsection{The parameter rebate and the capability-vs-scale curve shift}
\label{sec:framework-rebate}

Combining \ref{eq:rho}, \ref{eq:Cexplicit}, and Assumption~A1, the \textbf{parameter rebate} that morphology-aligned tokenization delivers for language $L$, namely the reduction in parameter count required to reach a fixed capability threshold, is:

\begin{equation}
\Delta N(L) = \frac{C_{\mathrm{implicit}}(L) - C_{\mathrm{explicit}}(L)}{\kappa} \approx \alpha \cdot \rho(L) \cdot H(L),
\label{eq:rebate}
\end{equation}

where $\kappa$ is the bits-per-parameter coefficient of the underlying architecture. The language-invariant constants are absorbed into a single $\alpha$, and the framework's prediction collapses to a single expression: the rebate is proportional to a language's effective grammatical information content.

Equation~\ref{eq:rebate} carries a matching prediction at the capability level. Let $N_{\mathrm{capability}}(T, L, \mathrm{tok})$ denote the parameter count at which a transformer trained on language $L$ with tokenization regime \textit{tok} first reaches capability target $T$, where $T$ can be defined on any continuous or discrete metric (a loss threshold, a task score, accuracy on a held-out probe). The framework predicts:

\begin{equation}
N_{\mathrm{capability}}(T, L, \mathrm{morph}) = N_{\mathrm{capability}}(T, L, \mathrm{BPE}) - \alpha \cdot \rho(L) \cdot H(L).
\label{eq:curveshift}
\end{equation}

Equation~\ref{eq:curveshift} is neutral between the two live views of how capability scales with parameter count. Under the sharp-emergence reading \citep{wei2022}, it shifts the threshold at which a capability first appears. Under the smooth-improvement reading \citep{schaeffer2023}, it shifts the entire continuous capability-vs-scale curve to the left by $\Delta N(L)$. Both readings predict the same directional effect, and differ only in the shape of the curve being shifted. The framework does not need to choose between them.

Equation~\ref{eq:curveshift} is also the formal statement of the \emph{edge-hostable reasoning} motivation discussed in \S\ref{sec:intro}. Let $N_{\mathrm{edge}}$ denote the parameter budget of a target edge profile, and $T_{\mathrm{reason}}$ a capability threshold associated with reasoning. The condition for edge-hostable reasoning in language $L$ under morphology-aligned tokenization is:

\begin{equation}
N_{\mathrm{capability}}(T_{\mathrm{reason}}, L, \mathrm{morph}) \le N_{\mathrm{edge}}.
\label{eq:edgecondition}
\end{equation}

Whether any given language-capability pair satisfies~\ref{eq:edgecondition} depends on $\alpha$, that is, on the quantitative strength of the rebate, which~\ref{eq:rebate} and~\ref{eq:curveshift} leave as the empirical unknown the framework hands to the experiment in \S\ref{sec:results}. The framework supplies the functional form, the computable inputs ($H(L)$ and $\rho(L)$), and the interpretive scaffolding. What it does not supply, and cannot supply from structural considerations alone, is the value of $\alpha$.

\S\ref{sec:results} reports the result of estimating $\alpha$ from Phase~1 observations. The ordinal prediction of~\ref{eq:rebate}, that the rebate orders the three languages, is confirmed. The magnitude prediction, that a single $\alpha$ fits all three, is not. This pair of outcomes is the substantive empirical content of the paper, and we return to it under the name \textbf{Agglutinative Compounding Effect} in \S\ref{sec:results-naming}.

\begin{figure}[t]
\centering
\includegraphics[width=0.85\linewidth]{fig_03_curve_shift_schematic.pdf}
\caption{Predicted capability-vs-scale curve shift under morphology-aligned tokenization (Equation~\ref{eq:curveshift}). The framework is neutral between sharp-emergence \citep{wei2022} and smooth-improvement \citep{schaeffer2023} readings of the capability-vs-scale curve; only the magnitude of the shift $\Delta N(L)$ is predicted.}
\label{fig:curve-shift}
\end{figure}
